% !TITLE 雁门太守行
% !AUTHOR 李贺
% !DYNASTY 唐
% !GENRE 乐府诗
% !ORDER 53

\begin{poem}
黑云压城城欲摧\textsuperscript{1}，甲光向日金鳞开\textsuperscript{2}。
角声满天秋色里，塞上燕脂凝夜紫\textsuperscript{3}。
半卷红旗临易水\textsuperscript{4}，霜重鼓寒声不起\textsuperscript{5}。
报君黄金台上意\textsuperscript{6}，提携玉龙\textsuperscript{7}为君死。
\end{poem}

\begin{annotations}
\notetext{1}{黑云压城城欲摧：黑云（比喻敌军）压城而来，城墙似乎都要被压垮了。“摧”字制造了窒息感。}
\notetext{2}{甲光向日金鳞开：阳光穿透云层照在铠甲上，像金色鱼鳞一样闪闪发光。黑暗与金光对比，压抑中闪现希望。}
\notetext{3}{燕脂凝夜紫：胭脂色的血迹在夜色中凝结成紫色。燕脂，即胭脂，指战场上的血迹。李贺用颜色写战争的残酷。}
\notetext{4}{易水：河名，在今河北。荆轲刺秦时在此告别燕太子丹，唱“风萧萧兮易水寒，壮士一去兮不复还”。暗示将士已知必死。}
\notetext{5}{霜重鼓寒声不起：霜太重，战鼓被冻得声音发不出来。写边塞严寒和战斗条件之艰苦。}
\notetext{6}{黄金台上意：燕昭王筑黄金台招贤纳士，象征君王对人才的礼遇。黄金台，相传在易水附近。}
\notetext{7}{玉龙：宝剑的美称。提携玉龙——手握宝剑，为报君恩视死如归。}
\end{annotations}

\begin{appreciation}
这是李贺最出名的战争诗之一，也是唐诗里颜色用得最凶的一首。黑、金、红、紫，像用颜料泼出来的。

黑云压城城欲摧——黑云不是云，是敌军。压和摧连用，一座城要被压塌了，开篇就让人喘不上气。甲光向日金鳞开——一道光从云缝里漏下来，照在铠甲上，像金色的鱼鳞。黑暗里裂开一道金，守军还在。

角声满天秋色里，塞上燕脂凝夜紫——燕脂就是胭脂，这里说的是血。李贺不说血，说燕脂；不说红，说夜里凝成的紫。这是李贺最狠的地方：把残酷写到接近好看。读的时候觉得画面美，反应过来才觉得冷。

半卷红旗临易水——旗只展开一半，是急行军。易水两个字，让人想起荆轲：风萧萧兮易水寒，壮士一去兮不复还。那是明知回不来也要去的故事。霜重鼓寒声不起——霜太重，鼓都冻哑了。战鼓是军队的喉咙，喉咙哑了，这一仗有多难打，不用多说。

报君黄金台上意，提携玉龙为君死——黄金台是燕昭王筑来招贤的台子，玉龙是宝剑。报答知遇之恩，赴死。放在前面层层压下来的黑云紫血之后，这两句不是豪言壮语，是一句交代。

还记得《无衣》吗？岂曰无衣？与子同袍——那些一起修戈矛、同仇敌忾的人，上了战场，就是这首诗里的人。《采薇》里那个“昔我往矣，杨柳依依”的士兵，要是回得来，大概也就是这样一个个送走了战友。

李贺二十七岁就死了，一辈子没上过战场，身体也不好。最浓烈的战争画面，来自一个最没有力气的人——想象力能抵达人没去过的地方。

这首诗的颜色，希望你只在课本上见过。
\end{appreciation}
